\chapter{Phương pháp đề xuất}
\label{Chapter3}

\section{Phương pháp đề xuất} 

Dựa trên các nghiên cứu liên quan, nhóm đề xuất xây dựng một hệ thống học sâu đa phương thức nhằm khai thác đồng thời thông tin từ ảnh X-quang và dữ liệu lâm sàng văn bản. Đặc biệt, hệ thống được thiết kế theo hướng tích hợp tri thức y khoa nhằm nâng cao khả năng học ít mẫu và phát hiện dữ liệu ngoài phân phối. Cụ thể, kiến trúc bao gồm các thành phần chính như sau:
    \begin{itemize}
        \item \textbf{Bộ mã hóa ảnh:} Kế thừa lại mô hình ViT-Base của phương pháp BiomedCLIP \cite{BiomedCLIP}
        \item \textbf{Bộ mã hóa văn bản:} Kế thừa lại mô hình PubMedBERT của phương pháp BiomedCLIP \cite{BiomedCLIP}
        \item \textbf{Bộ dung hợp dữ liệu:} Nhóm hướng đến sử dụng cơ chế \textbf{\textit{Knowledge-guided Cross-Attention}} - Các khái niệm y học được trích xuất và liên kết từ văn bản lâm sàng thông qua sẽ được sử dụng để tạo thiên lệch (bias) trong quá trình tính toán cross attention. Cơ chế này giúp điều hướng mô hình tập trung vào các vùng ảnh có liên quan đến bệnh lý, thay vì các vùng nền không liên quan. Nhờ đó, mô hình không chỉ cải thiện hiệu năng mà còn cung cấp khả năng giải thích thông qua việc truy vết mối liên hệ giữa khái niệm y khoa và vùng ảnh được chú ý.
        \item \textbf{Bộ phân loại:} Sử dụng một mạng phân loại (Multi-Layer Perceptron) để dự đoán nhãn bệnh lý.
    \end{itemize}
    Đồng thời, các trọng số attention và thông tin từ tri thức y khoa được khai thác để trực quan hóa và giải thích kết quả dự đoán.

\section{Giai đoạn off-line}
Giai đoạn off-line nhằm xây dựng bộ tham số cho mô hình và thiết lập không gian biểu diễn chung giữa đặc trưng hình ảnh X-quang và ngữ nghĩa y khoa, dựa trên dữ liệu tĩnh gồm ảnh và báo cáo chẩn đoán.

\subsection{Các kỹ thuật được sử dụng}
\begin{itemize}
    \item \textbf{Tiền xử lý văn bản với UMLS}: Phương pháp sử dụng LLM để trích xuất các tình trạng y khoa cùng với mức độ và vị trí từ báo cáo chẩn đoán ảnh X-quang, sau đó ánh xạ chúng vào hệ thống phân loại UMLS (Unified Medical Language System) để chuẩn hóa và liên kết tri thức từ đó xây dựng được các câu mô tả nhằm làm giàu ngữ nghĩa.
    \item \textbf{Tinh chỉnh tham số của các bộ mã hóa bằng phương pháp LoRA (Low-Rank Adaptation)}: Thay vì tinh chỉnh toàn bộ trọng số của Vision Encoder và Text Encoder, phương pháp chỉ cập nhật một tập con nhỏ các tham số thích ứng hạng thấp (low-rank) được cấy song song vào các khối Attention. Điều này giúp giảm thiểu rủi ro quá khớp (overfitting) và bảo tồn tri thức nền tảng của các mô hình mã hóa.
    
    Để cập nhật trọng số của LoRA, nhóm sử dụng \textbf{Hàm mất mát Khớp ngữ nghĩa (Semantic Matching Loss - $\mathcal{L}_{SM}$)} nhằm đồng bộ hóa không gian biểu diễn đa phương thức. Thay vì tính toán trên các véc-tơ tổng, mô hình thiết lập ma trận tương đồng $s(I_i, T_j)$ dựa trên sự căn chỉnh giữa các vùng ảnh (patches) của ảnh $I_i$ và các từ vựng (tokens) của câu lệnh $T_j$. 

    Hàm $\mathcal{L}_{SML}$ được định nghĩa là trung bình cộng của hai hàm mất mát tương phản hai chiều (Bi-directional Contrastive Loss), bao gồm chiều từ ảnh sang chữ ($\mathcal{L}_{I2T}$) và từ chữ sang ảnh ($\mathcal{L}_{T2I}$):
    \begin{equation}
        \mathcal{L}_{I2T} = -\frac{1}{N} \sum_{i=1}^{N} \log \frac{\exp(s(I_i, T_i) / \tau)}{\sum_{j=1}^{N} \exp(s(I_i, T_j) / \tau)}
    \end{equation}

    \begin{equation}
        \mathcal{L}_{T2I} = -\frac{1}{N} \sum_{i=1}^{N} \log \frac{\exp(s(I_i, T_i) / \tau)}{\sum_{j=1}^{N} \exp(s(I_j, T_i) / \tau)}
    \end{equation}

    \begin{equation}
        \mathcal{L}_{SML} = \frac{1}{2} (\mathcal{L}_{I2T} + \mathcal{L}_{T2I})
    \end{equation}
    trong đó, $N$ là kích thước mini-batch, $s(I_i, T_j)$ là điểm tương đồng ngữ nghĩa tổng hợp giữa ảnh $i$ và văn bản $j$, và $\tau$ là siêu tham số nhiệt độ (temperature) giúp điều chỉnh độ nhạy của hàm phạt đối với các mẫu âm tính (negative samples). Bằng cách tối ưu $\mathcal{L}_{SML}$, mô hình ép các đặc trưng giải phẫu X-quang cục bộ phải hội tụ chặt chẽ về đúng định nghĩa y khoa của chúng trong không gian chung.

    \item \textbf{Dung hợp đặc trưng}: Sau khi tinh chỉnh, các đặc trưng từ Vision Encoder và Text Encoder được kết hợp thông qua cơ chế Cross-Attention có hướng dẫn tri thức. Cụ thể, các khái niệm y khoa được trích xuất từ văn bản sẽ được sử dụng để tạo thiên lệch (bias) trong quá trình tính toán attention, giúp mô hình tập trung vào các vùng ảnh có liên quan đến bệnh lý.
\end{itemize}

\subsection{Các ẩn số cần tìm}
Để tránh bùng nổ tài nguyên tính toán, hệ thống áp dụng kỹ thuật \textbf{Tinh chỉnh hiệu quả tham số thông qua LoRA (Low-Rank Adaptation)}. Các ẩn số cần tìm và cơ chế cập nhật trọng số (Weight Update) được định nghĩa như sau:

\begin{itemize}
    \item \textbf{Đóng băng trọng số gốc:} Toàn bộ trọng số của Vision Encoder ($W_{0}^{v}$) và Text Encoder ($W_{0}^{t}$) được giữ cố định. Gradient không được tính toán cho các lớp này, giúp bảo toàn tri thức nền tảng.
    \item \textbf{Cập nhật trọng số LoRA:} Các ẩn số cần tìm là các ma trận thích ứng hạng thấp (Low-rank matrices) $A$ và $B$ được cấy song song vào các khối Attention. Trọng số thực tế của mô hình được biểu diễn bằng công thức:
    \begin{equation}
        W = W_0 + \Delta W = W_0 + BA
    \end{equation}
    \item \textbf{Lan truyền ngược (Backpropagation):} Trong quá trình huấn luyện, đạo hàm từ hàm mất mát $\mathcal{L}_{Align}$ chỉ chảy qua và cập nhật ma trận $A \in \mathrm{R}^{r \times k}$ và $B \in \mathrm{R}^{d \times r}$ (với hạng $r \ll \min(d, k)$), cùng với trọng số của hai lớp chiếu (Projection Heads) là $W_{proj\_v}$ và $W_{proj\_t}$. Tổng số tham số cần học (learnable parameters) chiếm dưới 5\% toàn mạng.
\end{itemize}

\subsection{Đóng góp của phương pháp}
\begin{itemize}
    \item \textbf{Triệt tiêu giằng xé đạo hàm (Gradient Conflict):} Việc gộp toàn bộ tri thức vào một hàm mất mát duy nhất ($\mathcal{L}_{Align}$) giúp luồng đạo hàm hội tụ ổn định, loại bỏ rủi ro xung đột khi phải tinh chỉnh nhiều siêu tham số ($\lambda$) trong các phương pháp cũ.
    \item \textbf{Chống sụp đổ ngữ nghĩa (Semantic Collapse):} Cơ chế cập nhật bằng LoRA giúp Text Encoder bảo vệ được không gian ngữ nghĩa y khoa nguyên bản (như PubMedBERT) khỏi việc bị bóp méo do quá khớp (overfitting) với các đặc trưng nhiễu của ảnh.
    \item \textbf{Tối ưu không gian Từ vựng mở (Open-vocabulary):} Phép chiếu token-patch trực tiếp từ hàm mất mát ép hệ thống phải định vị chính xác tổn thương cục bộ, tạo tiền đề toán học vững chắc cho bài toán Few-shot và phát hiện dị vật ngoài phân phối (OOD).
\end{itemize}

\section{Giai đoạn on-line}
\subsection{Các kỹ thuật được sử dụng}